\section{Experiment}

\subsection{Challenging Testbeds}

We conduct experiments on two challenging tasks, namely, commonsense machine translation (i.e., Common~MT), and counter-intuitive arithmetic reasoning (i.e., Counter-Intuitive AR), which require deep levels of contemplation for LLMs. Please refer to Appendix~\ref{app:testbeds} for more details.

\paragraph{Commonsense Machine Translation}
The Common~MT dataset is composed of Chinese$\Rightarrow$English translation examples~\citep{he-etal-2020-box}, which are used to examine three types of ambiguity resolution abilities of translation models, covering lexical
and contextless/contextual syntactic ambiguity. Within the challenging part of Common~MT, the authentic translation of each source sentence requires a proper understanding of common sense knowledge. While these ambiguous sentences might appear to have a straightforward translation, such a literal interpretation is erroneous. Failure to address such ambiguities may result in inaccurate translations.

\paragraph{Counter-Intuitive Arithmetic Reasoning}
Previous studies on thinking hierarchy~\citep{kahneman2011thinking} suggest that we humans have a fast and intuitive system and a slow and logical system, and tend to run the lower level system before the higher level one. Inspired by this, we created a more challenging dataset named Counter-Intuitive Arithmetic Reasoning (CIAR) to evaluate the reasoning abilities of LLMs at deep levels. 
Our Counter-Intuitive AR dataset contains 200 questions collected from elicitation 
questions~\citep{kong2022eliciting}\footnote{https://elicitation.info/questionnaire/1/}, 
web data\footnote{https://www.geeksforgeeks.org/puzzles/} and additional manual derivatives of these questions.
Compared to the commonly-used datasets, e.g., MultiArith~\citep{roy2015solving}, GSM8K~\citep{cobbe2021training},  our dataset presents two distinct challenges:
\begin{itemize}[leftmargin=10pt]
    \item \textit{Resistance to Intuition.} 
    The questions are embedded in hidden traps designed to elicit intuitive and appealing answers that are often incorrect. This feature evaluates the abilities of LLMs to resist the traps of superficial expressions.
    \item \textit{Multi-Step Reasoning.} 
    Each correct answer within the dataset requires a rigorous multi-step reasoning process, thereby evaluating the capacity of LLMs to engage in complex decision-making and problem-solving.
\end{itemize}

\subsection{Setups}
\paragraph{Input Format.}
Our experiments are performed in zero-shot instructions (setting temperature to 0).
For all used datasets, we use a unified prompt to make LLMs give explanations and answers.
We present the inputs to agents through $<$debate topic$>$ as mentioned in Section~\ref{sec:method}. For example, if we want to translate ``\zh{吃掉敌人一个师}'' from Chinese to English, we will set the $<$debate topic$>$ as \textit{``What is the correct English translation of the following Chinese text: \zh{吃掉敌人一个师}''}. For QA task, we employ the same prompt except set the $<$debate topic$>$ to the arithmetic question.

\paragraph{Backbone Models.}
In this work, we mainly use three agents in our MAD framework, including two debaters (i.e., affirmative and negative) and a judge. 
We assess two open-source (i.e., \texttt{vicuna-7b-v1.5-16k}\footnote{https://huggingface.co/lmsys/vicuna-7b-v1.5-16k} and \texttt{vicuna-13b-v1.5-16k}\footnote{https://huggingface.co/lmsys/vicuna-13b-v1.5-16k}) and two api-based LLMs (i.e., \texttt{GPT-3.5-Turbo-0301} and \texttt{GPT-4-0314}).


\begin{table*}[t]
\centering
\resizebox{\linewidth}{!}{
\begin{tabular}{l ccc ccc ccc}
\toprule
\multirow{2}{*}{\bf Method} & \multicolumn{3}{c}{\bf \normalsize Lexical} & \multicolumn{3}{c}{\bf \normalsize Contextless} & \multicolumn{3}{c}{\bf \normalsize Contextual} \\
\cmidrule(lr){2-4} \cmidrule(lr){5-7} \cmidrule(lr){8-10}
 & \small COMET & \small BLEURT & \small HUMAN & \small COMET & \small BLEURT & \small HUMAN & \small COMET & \small BLEURT & \small HUMAN\\
\midrule
\bf GPT-4  & \textbf{82.0} & 70.1 & 3.41 & 84.7 & 73.6 & 3.63 & 85.0 & 73.7 & 3.65 \\
\midrule
\bf \texttt{Turbo}  & 80.3 & 68.2 & 3.14 & 84.0 & 72.9 & 3.43 & 84.9 & 73.4 & 3.57 \\
~+ \textbf{Rerank}  & 80.9 & 68.6 & 3.16 & 84.5 & 73.2 & 3.46 & \textbf{85.3} & 73.9 & 3.58 \\
~+ \textbf{MAPS}  & 81.9 & 70.1 & 3.43 & 84.2 & 73.5 & 3.45 & 85.2 & \textbf{74.0} & 3.56 \\
~+ \textbf{\small Self-Reflect}  & 81.0 & 69.1 & 3.43 & 83.6 & 72.2 & 3.46 & 84.9 & 73.5 & 3.63 \\
\hdashline
~+ \textbf{MAD}  & \textbf{82.0} & \textbf{70.9} & \textbf{3.78} & \textbf{84.8} & \textbf{73.7} & \textbf{3.67} & \textbf{85.3} & \textbf{74.0} & \textbf{3.67} \\
\midrule
\bf \texttt{Vicuna-7b}  & 74.9 & 62.0 & 2.55  & 78.3 & 64.6 &  2.53  & 80.2 & 68.2 &  3.23  \\
~+ \textbf{MAD}  & 75.6 & 62.6 & 2.67  & 78.6 & 66.0 &  2.69  & 81.8 & 69.9 &  3.27  \\
\bf \texttt{Vicuna-13b}  & 76.6 & 63.7 & 2.81  & 77.6 & 66.8 &  3.04  & 82.2 & 70.0 &  3.37  \\
~+ \textbf{MAD}  & 77.2 & 65.1 & 2.96  & 80.1 & 67.3 & 3.11   & 82.6 & 70.9 &  3.45  \\
\bottomrule
\end{tabular}
}
\vspace{-5pt}
\caption{Translation performance on Common~MT. Note that Rerank and MAPS use the external quality estimation tool to select the best translation from multiple translation candidates. HUMAN: direct assessment of translation quality from human evaluators on a scale ranging from 1 to 5.}
\label{tab:common-mt}
\vspace{-10pt}
\end{table*}

\paragraph{Compared Methods.}
Generally, we compare our MAD framework with baseline models and Self-Reflect on both tasks. We also include other baseline methods individually, namely, Rerank and MAPS for Common~MT, CoT and Self-Consistency for Counter-Intuitive~AR. Below elaborates the details of them:
\begin{itemize}[leftmargin=10pt]
    \item \textbf{Self-Reflect}~\citep{shinn2023reflexion}: This approach requires the LLM to refine its translation until it deems the current output satisfactory.
    
    \item \textbf{Rerank}~\citep{he2023exploring}: We sample the translations from the LLM for four times, from which we select the best candidate based on a quality estimation (QE) HUMANr\footnote{We use \texttt{wmt21-comet-qe-da} as the QE HUMANr.}.
    This approach can be seen as analogous to self-consistency~\citep{wang2022self}, where the majority voting is replaced by an external QE HUMANr. 
    \item \textbf{MAPS}~\citep{he2023exploring}: This method enables LLMs to mimic the human translation process: analyze before translate, which can be viewed as a chain-of-thought method applied to translation. 
    
    \item \textbf{CoT}~\citep{kojimalarge}: This approach concatenates a trigger sentence ``Let’s think step by step''  to the test question.
    \item \textbf{Self-Consistency}~\citep{wang2022self}: This method samples multiple responses and determines the final answer through a majority vote. 
   
\end{itemize}

All agents in our experimental setup, such as debaters and judge, are large language models. Here, we implement the methods on top of \texttt{GPT-3.5-Turbo} and \texttt{Vicuna} models.


\paragraph{Evaluation Metrics.}
For Counter-Intuitive AR, we report the accuracy~(ACC) of predictions.
For Common~MT, we adopt automatic metrics like COMET\footnote{\texttt{\url{https://github.com/Unbabel/COMET/}, Unbabel/wmt22-comet-da}} and BLEURT\footnote{\texttt{\url{https://github.com/google-research/bleurt}, BLEURT-20}}, which are widely adopted evaluation metrics for LLM-based translation literature~\citep{he2023exploring,gpt-mt-2023,garcia2023unreasonable,pilault2023interactive}. 
In addition, we also employ professional human translators to directly assess the translation results, measuring translation quality on a scale ranging from 1 to 5.

\subsection{Results on Common~MT}
\label{subsec_mt}

\paragraph{Results.}
In Common~MT test set, we focus more on the translation accuracy of specific words and whether they conform to common sense. However, such minor variations at token level are difficult to reflect on automatic metrics. We therefore provide human HUMAN to evaluate these methods more accurately.
Table~\ref{tab:common-mt} presents the experimental results. MAPS and Self-Reflec achieve improvements over baseline \texttt{GPT-3.5-Turbo}. Remarkably, our proposed MAD, by utilizing GPT-3.5 as the backbone model, has demonstrated significant advancements over GPT-4 across both automatic and human evaluation metrics. 


\begin{table}[!tb]
\centering
\resizebox{\columnwidth}{!}{
\begin{tabular}{l p{5cm}}
\toprule
\bf Source & \zh{\underline{吃掉}敌人一个师。}\\
\bf Correct Ref. & \textcolor{blue}{Destroy} a division of the enemy.\\
\bf Incorrect Ref. & \textcolor{red}{Eat up} an enemy division.\\
\midrule
\bf \texttt{GPT-4} & \textcolor{red}{Eat up} an enemy division. \\
\bf \texttt{GPT-3.5-Turbo} & \textcolor{red}{Eat up} an enemy division. \\
~~+ \textbf{Self-Reflect}  & \textcolor{red}{Eat up} an enemy division. \\
\hdashline
~~+ \textbf{MAD} & \textcolor{blue}{Eliminate} an enemy division.  \\
\bottomrule
\end{tabular}
}
\vspace{-5pt}
\caption{Example translations generated by different methods. Best viewed in color.}
\label{tab:mt-case}
\vspace{-10pt}
\end{table}

\paragraph{Case Study.} 
Table~\ref{tab:mt-case} shows example translations generated by baseline \texttt{GPT-3.5-Turbo} and the proposed MAD.  We can find that the baseline \texttt{GPT-3.5-Turbo} (even the more powerful \texttt{GPT-4}) incorrectly translates the source words literally. Because of the DoT issue, Self-Reflect cannot rectify the literal translation. The proposed MAD framework, which explores divergent chain-of-thoughts, can generate the free translation of the underlined words within the source sentences. 

\subsection{Results on Counter-Intuitive AR}
\label{subsec_ar}

\begin{table}[t!]
\centering
\begin{tabular}{l c}
\toprule
\bf Method  & \multicolumn{1}{c}{\bf ACC (\%)} \\
\midrule
\bf\texttt{GPT-4}         &  51.0 \\
\hline
\bf\texttt{GPT-3.5-Turbo} & 26.0 \\
\bf ~~+ CoT        & 28.0       \\
\bf ~~+ Self-Consistency      & 29.5   \\
\bf ~~+ Self-Reflect  &  27.5         \\ 
\hdashline
\bf ~~+ MAD               & \textbf{37.0} \\
\bottomrule
\end{tabular}
\vspace{-5pt}
\caption{Accuracy on Counter-Intuitive AR.}
\label{tab:CIAR}
\vspace{-10pt}
\end{table}

\paragraph{Results.} Table~\ref{tab:CIAR} lists the results in terms of reasoning accuracy.
We can observe that Self-Reflect only marginally improves over the baseline \texttt{GPT-3.5-Turbo}, while CoT and Self-Consistency bring more improvements.
Our MAD framework, though not as good as \texttt{GPT-4}, outperforms all the other compared methods based on \texttt{GPT-3.5-Turbo}, which further demonstrates its effectiveness. We also validate MAD on math and symbolic reasoning tasks and report our results in Appendix~\ref{app:math_symbolic_results}.

\paragraph{Case Study.} Figure~\ref{fig:method} shows an example on Counter-Intuitive AR. We find both CoT and Self-Reflect fail to reach the right answer by mistakenly outputing $3$. With divergent thinking, our MAD framework emerges ``\textit{we need to consider both the rotation around circle B and the rotation of circle A itself}'' and find the correct answer $4$.
